% Input and encoding
\usepackage[utf8]{inputenc}
\usepackage[utf8]{vietnam}
\usepackage{lmodern}

% Math packages
\usepackage{amsmath,amssymb,amsthm}
\usepackage{mathtools}
\usepackage{physics} % For better physics notation

% Common math symbols
\newcommand{\N}{\mathbb{N}}
\newcommand{\Z}{\mathbb{Z}}
\newcommand{\Q}{\mathbb{Q}}
\newcommand{\R}{\mathbb{R}}
\newcommand{\C}{\mathbb{C}}
\newcommand{\F}{\mathbb{F}}

% Graphics and colors
\usepackage{graphicx}
\usepackage[dvipsnames]{xcolor}
\usepackage{tikz,tikz-3dplot}
\usepackage{pgfplots}
\pgfplotsset{compat=1.18}

% TikZ libraries
\usetikzlibrary{
    calc,trees,positioning,arrows,chains,shapes.geometric,
    decorations.pathreplacing,decorations.pathmorphing,shapes,
    matrix,shapes.symbols,fit,arrows.meta,backgrounds,
    decorations.markings,patterns,3d
}

% Tables and captions
\usepackage{booktabs}
\usepackage{array}
\usepackage{caption}
\usepackage{subcaption}
\setbeamertemplate{caption}[numbered]

% Text formatting and algorithms
\usepackage{ulem}
\usepackage{algorithm}
\usepackage{algorithmic}
% \usepackage{algpseudocode}
\usepackage{listings}
% \usepackage{enumitem}

% Hyperlinks
\usepackage{hyperref}
\hypersetup{
    colorlinks=true,
    linkcolor=blue,
    filecolor=magenta,
    urlcolor=cyan,
    citecolor=green
}

% =============================================================================
% COLOR SCHEME
% =============================================================================

\definecolor{primaryblue}{HTML}{0065BD}
\definecolor{secondaryblue}{HTML}{4A90E2}
\definecolor{accentorange}{HTML}{FF6B35}
\definecolor{darkgray}{HTML}{2C3E50}
\definecolor{lightgray}{HTML}{ECF0F1}
\definecolor{successgreen}{HTML}{27AE60}
\definecolor{warningred}{HTML}{E74C3C}

% =============================================================================
% BEAMER THEME CONFIGURATION
% =============================================================================

\mode<presentation>{
    \usetheme{default}
    \useoutertheme{infolines}
    \useinnertheme{circles}
    \setbeamercovered{transparent}
    \setbeamertemplate{theorems}[numbered]
    \setbeamertemplate{navigation symbols}{}
}

% Custom color theme
\setbeamercolor{structure}{fg=primaryblue}
\setbeamercolor{frametitle}{bg=primaryblue,fg=white}
\setbeamercolor{title}{fg=primaryblue}
\setbeamercolor{subtitle}{fg=darkgray}
\setbeamercolor{author}{fg=darkgray}
\setbeamercolor{block title}{bg=secondaryblue,fg=white}
\setbeamercolor{block body}{bg=lightgray,fg=darkgray}
\setbeamercolor{block title alerted}{bg=warningred,fg=white}
\setbeamercolor{block body alerted}{bg=warningred!10,fg=darkgray}
\setbeamercolor{block title example}{bg=successgreen,fg=white}
\setbeamercolor{block body example}{bg=successgreen!10,fg=darkgray}

% =============================================================================
% CUSTOM FRAME ENVIRONMENTS
% =============================================================================

\usepackage{framed}
\pgfmathsetseed{1}

% Define layers for parchment effect
\pgfdeclarelayer{background}
\pgfsetlayers{background,main}

% Parchment styles
\tikzset{
    normal border/.style={
        primaryblue!40!black!10, decorate,
        decoration={random steps, segment length=2.5cm, amplitude=.7mm}
    },
    torn border/.style={
        primaryblue!40!black!5, decorate,
        decoration={random steps, segment length=.5cm, amplitude=1.7mm}
    }
}

% Parchment frame macros
\def\parchmentframe#1{
    \tikz{
        \node[inner sep=2em] (A) {#1};
        \begin{pgfonlayer}{background}
            \fill[normal border] 
                (A.south east) -- (A.south west) -- 
                (A.north west) -- (A.north east) -- cycle;
        \end{pgfonlayer}
    }
}

\def\parchmentframetop#1{
    \tikz{
        \node[inner sep=2em] (A) {#1};
        \begin{pgfonlayer}{background}    
            \fill[normal border]
                (A.south east) -- (A.south west) -- 
                (A.north west) -- (A.north east) -- cycle;
            \fill[torn border]
                ($(A.south east)-(0,.2)$) -- ($(A.south west)-(0,.2)$) -- 
                ($(A.south west)+(0,.2)$) -- ($(A.south east)+(0,.2)$) -- cycle;
        \end{pgfonlayer}
    }
}

\def\parchmentframebottom#1{
    \tikz{
        \node[inner sep=2em] (A) {#1};
        \begin{pgfonlayer}{background}   
            \fill[normal border]
                (A.south east) -- (A.south west) -- 
                (A.north west) -- (A.north east) -- cycle;
            \fill[torn border]
                ($(A.north east)-(0,.2)$) -- ($(A.north west)-(0,.2)$) -- 
                ($(A.north west)+(0,.2)$) -- ($(A.north east)+(0,.2)$) -- cycle;
        \end{pgfonlayer}
    }
}

\def\parchmentframemiddle#1{
    \tikz{
        \node[inner sep=2em] (A) {#1};
        \begin{pgfonlayer}{background}   
            \fill[normal border]
                (A.south east) -- (A.south west) -- 
                (A.north west) -- (A.north east) -- cycle;
            \fill[torn border]
                ($(A.south east)-(0,.2)$) -- ($(A.south west)-(0,.2)$) -- 
                ($(A.south west)+(0,.2)$) -- ($(A.south east)+(0,.2)$) -- cycle;
            \fill[torn border]
                ($(A.north east)-(0,.2)$) -- ($(A.north west)-(0,.2)$) -- 
                ($(A.north west)+(0,.2)$) -- ($(A.north east)+(0,.2)$) -- cycle;
        \end{pgfonlayer}
    }
}

\newenvironment{parchment}[1][Example]{%
    \def\FrameCommand{\parchmentframe}%
    \def\FirstFrameCommand{\parchmentframetop}%
    \def\LastFrameCommand{\parchmentframebottom}%
    \def\MidFrameCommand{\parchmentframemiddle}%
    \vskip\baselineskip
    \MakeFramed{\FrameRestore}
    \noindent\tikz\node[inner sep=1ex, draw=white, fill=primaryblue!30, 
        anchor=west, overlay] at (0em, 2em) {\sffamily#1};\par
}{%
    \endMakeFramed
}

% =============================================================================
% TITLE PAGE CUSTOMIZATION
% =============================================================================

\makeatletter
\newcommand\titlegraphicii[1]{\def\inserttitlegraphicii{#1}}
\titlegraphicii{}

\setbeamertemplate{title page}{
    {\usebeamercolor[fg]{titlegraphic}\inserttitlegraphic\hfill\inserttitlegraphicii\par}
    \begin{centering}
        \vskip1em
        \begin{beamercolorbox}[sep=8pt,center]{title}
            \usebeamerfont{title}\inserttitle\par%
            \ifx\insertsubtitle\@empty%
            \else%
                \vskip0.5em%
                {\usebeamerfont{subtitle}\usebeamercolor[fg]{subtitle}\insertsubtitle\par}%
            \fi%     
        \end{beamercolorbox}%
        \vskip1em
        \begin{beamercolorbox}[sep=8pt,center]{author}
            \usebeamerfont{author}\insertauthor
        \end{beamercolorbox}
        \vskip1em
        \begin{beamercolorbox}[sep=8pt,center]{institute}
            \usebeamerfont{institute}\insertinstitute
        \end{beamercolorbox}
        \vskip1em
        \begin{beamercolorbox}[sep=8pt,center]{date}
            \usebeamerfont{date}\insertdate
        \end{beamercolorbox}
    \end{centering}
    \vfill
}

% =============================================================================
% HEADER AND FOOTER
% =============================================================================

\setbeamertemplate{frametitle}{%
    \nointerlineskip%
    \begin{beamercolorbox}[wd=\paperwidth,ht=2.5ex,dp=1.125ex]{frametitle}
        \hspace*{1em}\insertframetitle%
        \ifx\insertframesubtitle\@empty%
        \else%
            \par\hspace*{1em}{\usebeamerfont{framesubtitle}\insertframesubtitle}%
        \fi%
    \end{beamercolorbox}%
}

\setbeamertemplate{footline}{
    \leavevmode%
    \hbox{%
        \begin{beamercolorbox}[wd=.25\paperwidth,ht=2.25ex,dp=1ex,center]{author in head/foot}%
            \usebeamerfont{author in head/foot}\insertpresenter
        \end{beamercolorbox}%
        \begin{beamercolorbox}[wd=.5\paperwidth,ht=2.25ex,dp=1ex,center]{title in head/foot}%
            \usebeamerfont{title in head/foot}\insertshorttitle
        \end{beamercolorbox}%
        \begin{beamercolorbox}[wd=.25\paperwidth,ht=2.25ex,dp=1ex,center]{date in head/foot}%
            \insertframenumber{} / \inserttotalframenumber\hspace*{1ex}
        \end{beamercolorbox}%
    }%
    \vskip0pt%
}

% =============================================================================
% SECTION PAGES
% =============================================================================

\setbeamertemplate{section page}{
    \begin{centering}
        \vskip2em
        \begin{beamercolorbox}[sep=12pt,center]{part title}
            \usebeamerfont{section title}\insertsection\par
        \end{beamercolorbox}
        \vskip2em
    \end{centering}
}

% \AtBeginSection[]{
%     \begin{frame}[plain,noframenumbering]
%         \sectionpage
%         \tableofcontents[currentsection,hideothersubsections]
%     \end{frame}
% }

% =============================================================================
% THEOREM ENVIRONMENTS
% =============================================================================

\usepackage{thmtools}
\usepackage[framemethod=TikZ]{mdframed}
\mdfsetup{skipabove=1em,skipbelow=1em}

% Theorem styles
\declaretheoremstyle[
    headfont=\bfseries\color{successgreen!70!black},
    bodyfont=\normalfont,
    % mdframed={
    %     linewidth=2pt,
    %     rightline=false, topline=false, bottomline=false,
    %     linecolor=successgreen, backgroundcolor=successgreen!5,
    % }
]{thmwhitebox}

\declaretheoremstyle[
    headfont=\bfseries\color{successgreen!70!black},
    bodyfont=\normalfont,
    mdframed={
        linewidth=2pt,
        rightline=false, topline=false, bottomline=false,
        linecolor=successgreen, backgroundcolor=successgreen!5,
    }
]{thmgreenbox}

\declaretheoremstyle[
    headfont=\bfseries\color{primaryblue!70!black},
    bodyfont=\normalfont,
    mdframed={
        linewidth=2pt,
        rightline=false, topline=false, bottomline=false,
        linecolor=primaryblue, backgroundcolor=primaryblue!5,
    }
]{thmbluebox}

\declaretheoremstyle[
    headfont=\bfseries\color{primaryblue!70!black},
    bodyfont=\normalfont,
    mdframed={
        linewidth=2pt,
        rightline=false, topline=false, bottomline=false,
        linecolor=primaryblue
    }
]{thmblueline}

\declaretheoremstyle[
    headfont=\bfseries\color{warningred!70!black},
    bodyfont=\normalfont,
    mdframed={
        linewidth=2pt,
        rightline=false, topline=false, bottomline=false,
        linecolor=warningred, backgroundcolor=warningred!5,
    }
]{thmredbox}

\declaretheoremstyle[
    headfont=\bfseries\color{accentorange!70!black},
    bodyfont=\normalfont,
    mdframed={
        linewidth=2pt,
        rightline=false, topline=false, bottomline=false,
        linecolor=accentorange, backgroundcolor=accentorange!5,
    }
]{thmorangebox}

\declaretheoremstyle[
    headfont=\bfseries\color{warningred!70!black},
    bodyfont=\normalfont,
    numbered=no,
    mdframed={
        linewidth=2pt,
        rightline=false, topline=false, bottomline=false,
        linecolor=warningred, backgroundcolor=warningred!1,
    },
    qed=\qedsymbol
]{thmproofbox}

% Define theorem environments
\declaretheorem[style=thmwhitebox, name=Bài toán]{problem}
\declaretheorem[style=thmwhitebox, name=Định nghĩa]{definition}
\declaretheorem[style=thmwhitebox, name=Ví dụ]{example}
\declaretheorem[style=thmwhitebox, name=Nhận xét]{remark}
\declaretheorem[style=thmwhitebox, name=Định lý]{theorem}
\declaretheorem[style=thmwhitebox, name=Bổ đề]{lemma}
\declaretheorem[style=thmwhitebox, name=Mệnh đề]{proposition}
\declaretheorem[style=thmwhitebox, name=Hệ quả]{corollary}
\declaretheorem[style=thmwhitebox, name=Khẳng định]{claim}
\declaretheorem[style=thmwhitebox, name=Chú ý]{mynote}
\declaretheorem[style=thmwhitebox, name=Câu hỏi]{question}

% Proof environment
\declaretheorem[style=thmproofbox, name=Proof]{replacementproof}
\renewenvironment{proof}[1][\proofname]{\vspace{-10pt}\begin{replacementproof}}{\end{replacementproof}}

% =============================================================================
% CODE LISTINGS
% =============================================================================

\lstdefinestyle{moderncode}{
    basicstyle=\ttfamily\footnotesize,
    keywordstyle=\color{primaryblue}\bfseries,
    commentstyle=\color{successgreen}\itshape,
    stringstyle=\color{warningred},
    numberstyle=\tiny\color{darkgray},
    breaklines=true,
    breakatwhitespace=false,
    showspaces=false,
    showstringspaces=false,
    showtabs=false,
    tabsize=2,
    numbers=left,
    frame=leftline,
    frameround=fttt,
    framexleftmargin=2em,
    xleftmargin=2em,
    backgroundcolor=\color{lightgray!30},
    rulecolor=\color{primaryblue},
}

\lstset{style=moderncode}

% =============================================================================
% FLOWCHART SETTINGS
% =============================================================================

\tikzset{
    >=stealth',
    punktchain/.style={
        rectangle, 
        rounded corners=3pt, 
        draw=primaryblue, 
        very thick,
        fill=primaryblue!10,
        text width=6em,
        minimum height=2.5em, 
        text centered, 
        on chain
    },
    largepunktchain/.style={
        rectangle,
        rounded corners=3pt,
        draw=primaryblue, 
        very thick,
        fill=primaryblue!10,
        text width=10em,
        minimum height=2.5em,
        text centered,
        on chain
    },
    decision/.style={
        diamond,
        draw=warningred,
        very thick,
        fill=warningred!10,
        text width=4em,
        minimum height=2em,
        text centered,
        on chain
    },
    line/.style={draw, thick, ->},
    every join/.style={->, thick, shorten >=1pt},
    font={\fontsize{10pt}{12}\selectfont},
}

% =============================================================================
% UTILITY COMMANDS
% =============================================================================

% Highlighting
\newcommand{\highlight}[1]{\colorbox{accentorange!20}{#1}}
\renewcommand{\alert}[1]{\textcolor{warningred}{\textbf{#1}}}
\newcommand{\emphasis}[1]{\textcolor{primaryblue}{\textbf{#1}}}

% Math operators (check if already defined)
% \providecommand{\rank}{}
% \renewcommand{\rank}{\operatorname{rank}}
% \providecommand{\trace}{}
% \renewcommand{\trace}{\operatorname{tr}}
% \DeclareMathOperator{\diag}{diag}
% \providecommand{\span}{}
% \renewcommand{\span}{\operatorname{span}}
% \DeclareMathOperator{\dimension}{dim}
% \DeclareMathOperator{\ker}{ker}
% \DeclareMathOperator{\im}{im}


% =============================================================================
% CUSTOM TITLE PAGE
% =============================================================================
% Define custom commands for title page information
\newcommand{\university}[1]{\def\insertuniversity{#1}}
\newcommand{\school}[1]{\def\insertschool{#1}}
\newcommand{\faculty}[1]{\def\insertfaculty{#1}}
\newcommand{\department}[1]{\def\insertdepartment{#1}}
\newcommand{\presenter}[1]{\def\insertpresenter{#1}}
\newcommand{\president}[1]{\def\insertpresident{#1}}
\newcommand{\examinerone}[1]{\def\insertexaminerone{#1}}
\newcommand{\examinertwo}[1]{\def\insertexaminertwo{#1}}
\newcommand{\examinerthree}[1]{\def\insertexaminerthree{#1}}
\newcommand{\examinerfour}[1]{\def\insertexaminerfour{#1}}
\newcommand{\supervisorone}[1]{\def\insertsupervisorone{#1}}
\newcommand{\supervisortwo}[1]{\def\insertsupervisortwo{#1}}

\makeatletter
\setbeamertemplate{title page}{
    % \vbox{}
    \vfill
    \begingroup
        \centering
        % School name at top
        \begin{beamercolorbox}[sep=8pt,center]{institute}
            \usebeamerfont{institute}\large\textbf{\MakeUppercase\insertuniversity}\\
            \usebeamerfont{institute}\normalsize\MakeUppercase\insertschool\\
            \usebeamerfont{institute}\normalsize\MakeUppercase\insertfaculty
            % \usebeamerfont{institute}\normalsize\MakeUppercase\insertdepartment
        \end{beamercolorbox}\par
        
        \vskip1.em\par
        
        % Title
        \begin{beamercolorbox}[sep=8pt,center]{title}
            \usebeamerfont{title}\Large\textbf{\MakeUppercase\inserttitle}\par%
            \ifx\insertsubtitle\@empty%
            \else%
                \vskip0.5em%
                {\usebeamerfont{subtitle}\large\insertsubtitle\par}%
            \fi%     
        \end{beamercolorbox}\par
        
        \vskip1.em\par
        
        % Presenter
        \begin{beamercolorbox}[sep=8pt,center]{author}
            \usebeamerfont{author}
            % \begin{tabular}{@{}p{.35\textwidth}@{\hspace{3pt}}p{.5\textwidth}@{}}
            %     \textbf{Presenter:} \insertpresenter & \textbf{Supervisor 1:} \insertsupervisorone \\
            %     & \textbf{Supervisor 2:} \insertsupervisortwo \\
            % \end{tabular}
            % \begin{center}
            %     \textbf{Học viên cao học:} \insertpresenter \\ 
            %     \textbf{Người hướng dẫn khoa học:} \insertsupervisorone
            % \end{center}
            \begin{center}
                \textbf{}\insertpresenter \\ 
                % \textbf{Người hướng dẫn khoa học:} \insertsupervisorone
            \end{center}
        \end{beamercolorbox}
        
        \vskip1em\par
        
        % Date
        \begin{beamercolorbox}[sep=8pt,center]{date}
            \usebeamerfont{date}\insertdate
        \end{beamercolorbox}\par
        
    \endgroup
    \vfill
}
\makeatother